\documentclass[10pt,a4paper]{amsart}
\usepackage[margin=19mm]{geometry}
\usepackage{amsmath,amssymb,mathtools}
\usepackage[scr=rsfs,cal=euler]{mathalfa}
\usepackage{xfrac}
\usepackage[unicode,hidelinks,bookmarks=false]{hyperref}
\newcommand{\ie}{i.e.\ }
\newcommand{\cf}{cf.\ }
\newcommand{\resp}{respectively\ }
\newcommand{\Subsection}{\subsection}
\providecommand{\nobreakdash}{\nobreak\hbox{-}}
\setlength{\emergencystretch}{2em}
\multlinegap0pt
\usepackage{bm}
\usepackage{bbm}
\usepackage{dsfont}
\usepackage{xspace}
\usepackage{cancel}
\usepackage{mathtools}
\usepackage{amsthm}
\makeatletter
\@ifundefined{theorem}{\newtheorem{theorem}{Theorem}}{}
\@ifundefined{lemma}{\newtheorem{lemma}[theorem]{Lemma}}{}
\makeatother
\newcommand{\dt}{\tilde}
\title[Well-posedness and Stability Analysis of Suspension Bridge M]{Well-posedness and Stability Analysis of Suspension Bridge Models Coupled with Cattaneo Heat Conduction: The Role of Viscoelastic Memory}
\author{Jun Zhou, Jiamin He}
\date{Source: arXiv:2606.12786v1 (CC BY 4.0)}
\begin{document}
\maketitle

\noindent\emph{Source and licence.} Well-posedness and Stability Analysis of Suspension Bridge Models Coupled with Cattaneo Heat Conduction: The Role of Viscoelastic Memory. Version arXiv:2606.12786v1 (CC BY 4.0), distributed under the Creative Commons Attribution 4.0 International licence, as recorded in the arXiv licence metadata for this exact version and shown on the abstract page https://arxiv.org/abs/2606.12786v1. Full rights notice: licenses/arxiv-2606.12786-CC-BY-4.0.txt.\par\smallskip

\noindent\textit{Editorial scope.} Counted unit: the Lemma whose source label is \texttt{maximal-1} in the article (source0.tex lines 389--391, source label \texttt{maximal-1}), delivered together with the article's own complete original proof of it (source0.tex lines 392--575, one \texttt{proof} environment of 184 lines). No mathematics is written, repaired or re-derived: statement and proof are the article's own text line for line. Required context retained uncounted: definition:A-1 (lines 322--333). Every definition, display and label of those lines is retained unchanged. Omitted from this package and not redistributed: 1--321, 334--388, 576--1754. Nothing inside a retained excerpt is omitted or altered: every retained body is delivered exactly as the article's lines are written, so no omission receipt of any kind accompanies this package. Citations: the retained excerpts carry no citation command at all, so no bibliography entry of the article is transcribed and no reference is redistributed. Reference adaptation: none was needed and none was made; every reference inside the delivered unit resolves inside it, so no unresolved reference or citation is produced. Printed slips kept literal: no printed word, symbol, hypothesis or number of the counted statement or of its proof was altered. Layout and class: the article's own class is \texttt{article} with packages including diagbox, amsfonts, bbm, bm, pifont, mathptmx, amsmath, mathrsfs, hyperref, amssymb, color, amsthm, amstext, natbib; the wrapper is the lane's pinned \texttt{amsart} editorial wrapper at 10pt on A4, which restates the theorem-like environments the retained text needs and adds the article's own macro definitions that the retained text needs (\texttt{\textbackslash dt}), with extra wrapper packages (bm, bbm, dsfont, xspace, cancel, mathtools). The delivered numbering is the unit's own. Assets: no retained span contains a figure environment, a graphics-inclusion command, a PDF-page inclusion, a table or a TikZ picture; no image file is copied into this package, the package declares no asset dependency and no image file is redistributed with it. Licence: version arXiv:2606.12786v1 of this article is distributed under the Creative Commons Attribution 4.0 International licence, as recorded in the arXiv licence metadata for this exact version and shown on the abstract page https://arxiv.org/abs/2606.12786v1. A full rights notice is delivered at licenses/arxiv-2606.12786-CC-BY-4.0.txt. Characters: every retained excerpt is pure ASCII, so the delivered engine, XeLaTeX with its default Latin Modern text font, reports no missing character.
where \(\mathcal{A}_1: \mathcal{D}(\mathcal{A}_1) \subset \mathcal{H}_1 \to \mathcal{H}_1\) is a linear operator defined by \begin{equation}\label{definition:A-1}
	\mathcal{A}_1 U := \begin{pmatrix}
		v \\
		\rho^{-1} \left[ \alpha u_{xx} + \beta (\varphi - u) - \gamma v \right] \\
		\Phi \\
		\rho_1^{-1} \left[ \kappa (\varphi_x + \psi)_x - \beta (\varphi - u) - m \theta_x \right] \\
		\Psi \\
		\rho_2^{-1} \left[ b \psi_{xx} - \kappa (\varphi_x + \psi) + m\theta \right] \\
		\rho_3^{-1} \left[ -q_x - m (\Phi_x + \Psi) \right] \\
		\tau^{-1}(-\sigma q - \theta_x)
	\end{pmatrix}
\end{equation}

\begin{lemma}\label{maximal-1}
	$0 \in \varrho(\mathcal{A}_1)$ (the resolvent set of $\mathcal{A}_1$). Consequently, $\mathcal{A}_1$ is maximal.
\end{lemma}

\begin{proof}
	First, we will prove for any $F = (f_1,f_2,f_3,f_4,f_5,f_6,f_7,f_8) \in \mathcal{H}_1$, there exists a unique $$U = (u,v,\varphi,\Phi,\psi,\Psi,\theta,q)\in \mathcal{D}(\mathcal{A}_1)$$ such that \begin{align*}
		\mathcal{A}_1U = F.
	\end{align*}
	From the definition of $\mathcal{A}_1$ \eqref{definition:A-1}, the expression can be reformulated as the following equations
	\begin{equation}\label{eq:resolvent-system-A-1}
		\begin{cases}
			v = f_1, \\
			\alpha u_{xx} + \beta (\varphi - u) - \gamma v = \rho f_2, \\
			\Phi = f_3, \\
			\kappa (\varphi_x + \psi)_x - \beta (\varphi - u) - m \theta_x = \rho_1 f_4, \\
			\Psi = f_5, \\
			b \psi_{xx} - \kappa (\varphi_x + \psi) + m\theta = \rho_2 f_6, \\
			-q_x - m (\Phi_x + \Psi) = \rho_3 f_7, \\
			-\sigma q - \theta_x = \tau f_8.
		\end{cases}
	\end{equation}
	From equations \eqref{eq:resolvent-system-A-1}$_{1,3,5}$, it follows that
	\begin{equation}\label{df1}
		v = f_1 \in H_*^1, \quad \Phi = f_3 \in H_*^1, \quad \Psi = f_5 \in H_0^1.
	\end{equation}
	Equation \eqref{eq:resolvent-system-A-1}$_8$ yields
	$q_x=-\frac{1}{\sigma}(\theta_{xx}+\tau f_{8x}).$ By substituting the above equation into equation \eqref{eq:resolvent-system-A-1}$_7$, we obtain
	\begin{align*}
		- \theta_{xx} = - \sigma m(\Phi_x + \Psi) - \sigma \rho_3 f_7 + \tau f_{8x}.
	\end{align*}
	Then combined with equations \eqref{eq:resolvent-system-A-1}$_{2,4,6}$, we derive
	\begin{align}\label{eq:resolvent-system-1-A-1}
		\begin{cases}
			-\alpha u_{xx} - \beta(\varphi - u) = -\rho f_2 - \gamma v=-\rho f_2 - \gamma f_1 := g_1 \in L_*^2, \\
			-\kappa(\varphi_x + \psi)_x + \beta(\varphi - u) + m\theta_x = -\rho_1 f_4 := g_2 \in L_*^2, \\
			- b \psi_{xx} + \kappa(\varphi_x + \psi) - m\theta = -\rho_2 f_6  := g_3 \in L^2, \\
			- \theta_{xx} = - \sigma m(\Phi_x + \Psi) - \sigma \rho_3 f_7 + \tau f_{8x}=- \sigma m(f_{3x} + f_5) - \sigma \rho_3 f_7 + \tau f_{8x}:=g_4 \in H^{-1}.
		\end{cases}
	\end{align}
	The weak formulation of equations \eqref{eq:resolvent-system-1-A-1} is defined as follows
	\begin{align}\label{eq:resolvent-system-2-A-1}
		\begin{cases}
			\alpha(u_x, \tilde{u}_x) - \beta(\varphi - u, \tilde{u}) = (g_1, \tilde{u}), \\
			\kappa(\varphi_x + \psi, \dt \varphi_x) + \beta(\varphi - u, \dt \varphi) - m(\theta, \dt \varphi_x) = (g_2, \dt \varphi), \\
			b(\psi_x, \dt \psi_x) + \kappa(\varphi_x + \psi, \dt \psi) - m(\theta, \dt \psi) = (g_3, \dt \psi), \\
			\left(\theta_x, \frac{m^2 c_p}{\kappa } \dt \theta_x\right) = \left(g_4, \frac{m^2 c_p}{\kappa } \dt \theta\right),
		\end{cases}
	\end{align}
	for any $(\tilde{u}, \tilde{\varphi}, \tilde{\psi}, \tilde{\theta}) \in V := H_*^1 \times H_*^1 \times H_0^1 \times H_0^1$. Summing each equation in \eqref{eq:resolvent-system-2-A-1} yields the bilinear form $\mathcal{L}: V \times V \to \mathbb{R}$ and the linear form $\mathcal{F}:V \to \mathbb{R}$, satisfying
	\begin{equation}\label{LF}
		\mathcal{L}\left(
		\begin{pmatrix}u\\ \varphi \\ \psi \\ \theta\end{pmatrix},
		\begin{pmatrix}\tilde{u}\\ \tilde{\varphi} \\ \tilde{\psi} \\ \tilde{\theta}\end{pmatrix}
		\right) = \mathcal{F}\left(
		\begin{pmatrix}\tilde{u}\\ \tilde{\varphi} \\ \tilde{\psi} \\ \tilde{\theta}\end{pmatrix}
		\right).
	\end{equation}
	where $\mathcal{L}$ and $\mathcal{F}$ are defined as
	\begin{equation}\label{L}
		\begin{split}
			\mathcal{L}\left(
			\begin{pmatrix}u\\ \varphi \\ \psi \\ \theta\end{pmatrix},
			\begin{pmatrix}\tilde{u}\\ \tilde{\varphi} \\ \tilde{\psi} \\ \tilde{\theta}\end{pmatrix}
			\right) &= \alpha (u_x, \tilde{u}_x) + \beta (\varphi - u, \tilde{\varphi} - \tilde{u}) + \kappa (\varphi_x + \psi, \tilde{\varphi_x} + \tilde{\psi}) \\
			&\quad - m(\theta, \tilde{\varphi_x}+ \tilde{\psi}) + b (\psi_x, \tilde{\psi_x})   + \frac{m^2 c_p}{\kappa} (\theta_x, \tilde{\theta_x}),
		\end{split}
	\end{equation}
	\begin{align}\label{F}
		\mathcal{F}\left(
		\begin{pmatrix}\tilde{u}\\ \tilde{\varphi} \\ \tilde{\psi} \\ \tilde{\theta}\end{pmatrix}
		\right) &= \left(g_1, \tilde{u}\right) + \left(g_2, \tilde{\varphi}\right) + \left(g_3, \tilde{\psi}\right) + \left(g_4, \frac{m^2 c_p}{\kappa }\tilde{\theta}\right) \notag \\
		&=  - \gamma (f_1, \tilde{u}) - \rho (f_2, \tilde{u}) - \rho_1 (f_4, \tilde{\varphi}) - \rho_2 (f_6, \tilde{\psi})  \notag \\
		&\quad  - \frac{\sigma m^3 c_p }{\kappa } (f_{3x} + f_5, \tilde{\theta}) - \frac{\sigma \rho_3 m^2 c_p }{\kappa } (f_7, \tilde{\theta}) - \frac{\tau m^2 c_p}{\kappa}(f_8, \tilde{\theta_x}).
	\end{align}
	Clearly, $\mathcal{L}$ and $\mathcal{F}$ are bounded. We now prove that $\mathcal{L}$ is coercive. In fact, for any $(u, \varphi, \psi, \theta) \in V$, by Young's inequality and Poincar\'e's inequality, we get
	\begin{align*}\label{L-1}
		\mathcal{L}\left(
		\begin{pmatrix}u\\ \varphi \\ \psi \\ \theta\end{pmatrix},
		\begin{pmatrix}u\\ \varphi \\ \psi \\ \theta\end{pmatrix}
		\right) &= \alpha \| u_x \|^2 + \beta \| \varphi - u \|^2 + \kappa \| \varphi_x + \psi \|^2 + b \| \psi_x \|^2 + \frac{m^2 c_p}{\kappa} \| \theta_x \|^2 - m (\theta, \varphi_x + \psi)\\
		&\geq \alpha \| u_x \|^2 + \beta \| \varphi - u \|^2 + \kappa \| \varphi_x + \psi \|^2 + b \| \psi_x \|^2 + \frac{m^2 c_p}{\kappa} \| \theta_x \|^2 - |m|\|\theta\|\|\varphi_x+\psi\|\\
		&\geq \alpha \| u_x \|^2 + \beta \| \varphi - u \|^2 + \kappa \| \varphi_x + \psi \|^2 + b \| \psi_x \|^2 + \frac{m^2 c_p}{\kappa} \| \theta_x \|^2 - \frac{\kappa}{2}\|\varphi_x+\psi\|^2 - \frac{m^2}{2\kappa}\|\theta\|^2\\
		&\geq \alpha \| u_x \|^2 + \beta \| \varphi - u \|^2 + \kappa \| \varphi_x + \psi \|^2 + b \| \psi_x \|^2 + \frac{m^2 c_p}{\kappa} \| \theta_x \|^2  - \frac{\kappa}{2}\|\varphi_x+\psi\|^2 - \frac{m^2 c_p}{2\kappa}\|\theta_x\|^2\\	
		&= \alpha \| u_x \|^2 + \beta \| \varphi - u \|^2 + \frac{\kappa}{2} \| \varphi_x + \psi \|^2 + b \| \psi_x \|^2 + \frac{m^2 c_p}{2\kappa} \| \theta_x \|^2 \\
		&\geq c \left\|
		\begin{pmatrix}u\\ \varphi \\ \psi \\ \theta\end{pmatrix}
		\right\|_V^2, 	
	\end{align*}
	where $c$ is a positive constant. Therefore, $\mathcal{L}$ is coercive. By the Lax-Milgram Theorem, system \eqref{eq:resolvent-system-1-A-1} has a unique solution $(u, \varphi, \psi, \theta) \in V$.
	
From equation \eqref{eq:resolvent-system-A-1}$_8$, it follows that
	\begin{align}
		q=-\frac{1}{\sigma}(\theta_x+\tau f_8) \in L^2.
	\end{align}
	From \eqref{eq:resolvent-system-A-1}$_7$, we also obtain $q_x=-m(\Phi_x+\Psi)-\rho_3 f_7 \in L^2$. Thus $q \in H^1$.
	
	Using standard elliptic regularity theory for \eqref{eq:resolvent-system-1-A-1}$_{1,2,3}$, we derive
	\begin{align}\label{jlll}
		u, \varphi, \psi \in H^2.
	\end{align}
	
	Setting $\dt \varphi = \dt \psi = \dt \theta = 0$ in \eqref{LF} yields
	\begin{align*}
		\alpha(u_x,\dt u_x) = (\beta(\varphi - u) + g_1, \dt u), \quad \forall \dt u \in H_*^1.
	\end{align*}
	Since $u\in H^2$, we can integrate by parts and rearrange the terms to obtain
	\begin{equation}\label{vb}
		(\alpha u_{xx}+\beta(\varphi - u) + g_1,\dt u)=\alpha \left(u_x(L)\dt u(L)-u_x(0)\dt u(0)\right), \quad \forall \dt u \in H_*^1.
	\end{equation}
	
	For $m$ an integer with $m\gg1$, we choose a nonnegative smooth function $\xi_m(x)$ defined on $\left[0,\frac{1}{m}\right]$ such that
	\[
	\xi_m(0)=1, \ 0\le\xi_m(x)\le 1 \text{ for } x\in\left(0,\frac{1}{m}\right), \ \xi_m\left(\frac{1}{m}\right)=0.
	\]
	Define
	\[
	\tilde{u}_m(x)=
	\begin{cases}
		\xi_m(x), & 0\le x\le\frac{1}{m},\\
		0, & \frac{1}{m}\le x\le L-\frac{1}{m},\\
		-\xi_m(L-x), & L-\frac{1}{m}\le x\le L.
	\end{cases}
	\]
	A series of straightforward calculations shows that $\tilde{u}_m(x)\in H_*^1$ and $\|\tilde{u}_m\|^2 \leq \frac{2}{m}$. Then by \eqref{vb}, we have
	\begin{equation}\label{vb1}
		(\alpha u_{xx}+\beta(\varphi - u) + g_1,\tilde{u}_m)=\alpha \left(-u_x(L)-u_x(0)\right).
	\end{equation}
	By Cauchy-Schwarz's inequality, we derive
	\[
	\left| (\alpha u_{xx}+\beta(\varphi - u) + g_1,\tilde{u}_m)\right|\le\left(\frac{2}{m}\right)^{\frac{1}{2}}\|\alpha u_{xx}+\beta(\varphi - u) + g_1\|\to0
	\]
	as $m\to\infty$.
	 Letting $m\to\infty$ in \eqref{vb1} gives
	\begin{equation}\label{vb2}
		u_x(L)+u_x(0)=0.
	\end{equation}
	Let
	\[
	\tilde{u}_m(x)=
	\begin{cases}
		\xi_m(x), & 0\le x\le \frac{1}{m},\\
		0, & \frac{1}{m}\le x\le\frac{L}{2}-\frac{1}{m},\\
		-\xi_m(\frac{L}{2}-x), &
		\frac{L}{2}-\frac{1}{m}\le x\le\frac{L}{2},\\
		-\xi_m(-\frac{L}{2}+x), &
		\frac{L}{2}\le x\le\frac{L}{2}+\frac{1}{m},\\
		0, &
		\frac{L}{2}+\frac{1}{m}\le x\le L-\frac{1}{m},\\
		\xi_m(L-x), & L-\frac{1}{m}\le x\le L.
	\end{cases}
	\]
	A series of straightforward calculations shows that $\tilde{u}_m(x)\in H_*^1$ and $\|\tilde{u}_m\|^2 \le\frac{4}{m}$. Then by \eqref{vb}, we have
	\begin{equation}\label{gvb1}
		(\alpha u_{xx}+\beta(\varphi - u) + g_1,\tilde{u}_m)=\alpha \left(u_x(L)-u_x(0)\right).
	\end{equation}
	By Cauchy-Schwarz's inequality, we derive
	\[
	\left| (\alpha u_{xx}+\beta(\varphi - u) + g_1,\tilde{u}_m)\right|\le\frac{2}{m^{\frac{1}{2}}}\|\alpha u_{xx}+\beta(\varphi - u) + g_1\|\to0
	\]
	as $m\to\infty$.
    Letting $m\to\infty$ in \eqref{gvb1} gives
	\begin{equation}\label{gvb2}
		u_x(L)-u_x(0)=0.
	\end{equation}
	
	From \eqref{vb2} and \eqref{gvb2}, it follows that
	\[
	u_x(L)=u_x(0)=0.
	\]
	Substituting back into \eqref{vb}, we obtain
	\[
	(\alpha u_{xx}+\beta(\varphi - u) + g_1,\tilde{u})=0, \quad \forall \tilde{u}\in H_*^1,
	\]
	which implies
	\[
	\alpha u_{xx}+\beta(\varphi - u) + g_1=0.
	\]
	Thus, we have
	\begin{equation}\label{jl6}
		u\in H^2\text{ with } u_x\in H_0^1 \text{ satisfies } \alpha u_{xx}+\beta(\varphi - u) + g_1=0.
	\end{equation}
	
	By similar reasoning,
	\begin{equation}\label{jl7}
		\varphi\in H^2\text{ with } \varphi_x\in H_0^1\text{ satisfies } -\kappa(\varphi_x + \psi)_x + \beta(\varphi - u) + m\theta_x = g_2.
	\end{equation}
	Consequently, $U \in \mathcal{D}(\mathcal{A}_1)$ is obtained. Clearly, the above proof shows that $\|U\|_{\mathcal{H}_1}=\|\mathcal{A}_1^{-1}F\|_{\mathcal{H}_1} \leq C\|F\|_{\mathcal{H}_1}$ for some positive constant $C$. Thus, $0 \in \varrho(\mathcal{A}_1)$.
\end{proof}

\end{document}
